\usepackage{emoji}
\usepackage{hyperref}
\usepackage{cleveref}

\usepackage{soul}  % for the \hl command in snippets 
\usepackage{graphicx}

% table stuff
\usepackage{piton}
\usepackage{tabularray}
\UseTblrLibrary{tikz,diagbox,functional}
\usetikzlibrary{fadings}  % for the pqcrypto table
\usepackage{xcolor,luacolor}
\usepackage{tcolorbox}
\tcbuselibrary{breakable,skins}
\usepackage{lua-ul}
\LuaULSetHighLightColor{yellow!40}

\usepackage{adjustbox}  % for \adjustbox
\usepackage{changepage}  % for \adjustwidth
\usepackage{bbding}  % for \Checkmark and \XSolidBrush
\newcommand{\yes}{\textcolor{green!70!black}{\Checkmark}}
\newcommand{\no}{\textcolor{red!50}{\XSolidBrush}}

\usepackage{paralist}  % compactitem

\newtcolorbox{infobox}[1]{
    breakable,
    colback=green!15,
    colframe=green!65!black,
    fonttitle=\bfseries,
    title={#1},
    before skip balanced=\parskip,
    after skip balanced=\parskip,
    skin=enhanced,
}

\newtcolorbox{errorbox}[1]{
    breakable,
    colback=red!15,
    colframe=red!65!black,
    fonttitle=\bfseries,
    title={#1},
    before skip balanced=\parskip,
    after skip balanced=\parskip,
    skin=enhanced,
}

\newtcolorbox{questionbox}[1]{
    breakable,
    colback=violet!15,
    colframe=violet!65!black,
    fonttitle=\bfseries,
    title={#1},
    before skip balanced=\parskip,
    after skip balanced=\parskip,
    skin=enhanced,
}

\newtcolorbox{warningbox}[1]{
    breakable,
    colback=orange!15,
    colframe=orange!65!black,
    fonttitle=\bfseries,
    title={#1},
    before skip balanced=\parskip,
    after skip balanced=\parskip,
    skin=enhanced,
}

\newtcolorbox{standalonebox}{
    colback=white,
    hbox,
    colframe=white,
    left*=0mm,
    right*=0mm,
    bottom=0mm,
    top=0mm,
    boxsep=0mm,
    sharp corners,
    skin=enhanced,
}

% TODO find a way to avoid having to escape the special characters
\newcommand{\snippet}[1]{{\sethlcolor{gray!30}\hl{\texttt{#1}}}}


% Latex syntax highlighting with piton
\NewPitonLanguage{LaTeX}{keywordsprefix = \ , alsoother = @_ , moredirectives={&}, morecomment = [l]\%}
\PitonOptions{language = LaTeX, width=min, detected-commands = {highLight}, splittable=4}
\SetPitonStyle[LaTeX]{ Number = , Directive= \color{red!50}}

%%%%%%%%%%%
% Display the code, render it, and create a copy-pastable annotation on the right side
%%%%%%%%%%%

% Almost copy-pasted from https://github.com/fpantigny/nicematrix/blob/master/nicematrix.tex#L79
\ExplSyntaxOn
\dim_new:N \l__pantigny_width_dim
\bool_new:N \l__pantigny_show_output_bool
\bool_new:N \l__pantigny_small_bool
\bool_new:N \l__pantigny_execute_bool
\bool_new:N \l__pantigny_footnotesize_bool 

\keys_define:nn { pantigny }
  {
    width .dim_set:N = \l__pantigny_width_dim ,
    annotation .code:n = \PitonOptions { annotation = #1 } ,
    show-output .bool_set:N = \l__pantigny_show_output_bool ,
    show-output .default:n = true,
  }

\NewPitonEnvironment { Code } { O { } }
  {
    \begin{tcolorbox}[colback=blue!3, boxrule=0pt]
        \dim_zero:N \l__pantigny_width_dim
        \PitonOptions { annotation, indent-broken-lines, continuation-symbol=\empty, end-of-broken-line=\empty, continuation-symbol-on-indentation=\empty }
        \keys_set:nn { pantigny } { #1 }
        \group_begin:
        \char_set_catcode_other:N |
        \color{gray}

        \small
        \dim_compare:nNnT \l__pantigny_width_dim > \c_zero_dim 
          { 
            \PitonOptions { width = \l__pantigny_width_dim  } 
            \begin{minipage}[c]{\l__pantigny_width_dim} 
          }
    }
    {
        \dim_compare:nNnT \l__pantigny_width_dim > \c_zero_dim 
          { \end{minipage} }
        \group_end:
        \bool_if:NTF \l__pantigny_show_output_bool
          {
                % \smallskip
                \vspace{-0.5em}
                \begin{center}\begin{standalonebox}
                    \lua_now:n { tex.print(piton.get_last_code():explode(string.char(10)))}
                \end{standalonebox}
                \end{center}
          }
      	  {
      	  	 \vspace{-0.2em}
      	  }
        \end{tcolorbox}
    }
\ExplSyntaxOff

%%%%
% for drawing the table in appendix
%%%%

% Inspired from https://tex.stackexchange.com/a/750321/430417
\IgnoreSpacesOn

\tlNew \gTikzPercentagePathsTl


\prgNewFunction \colMaxValue {m} {  
	\fpSet \lTmpaFp {\fpEval {\cellGetText{2}{#1}}}
	\intStepOneInline {2} {\arabic{rowcount}} {
		\fpSet \lTmpaFp {\fpMathMax{\lTmpaFp}{\fpEval {\cellGetText{##1}{#1}}}}
	}
	\prgReturn {\fpUse\lTmpaFp}
}

\prgNewFunction \barPlotPaths {m} {
	\tlClear \gTikzPercentagePathsTl
	\fpSet \lTmpaFp {\colMaxValue{#1}}
	\intStepOneInline {2} {\arabic{rowcount}} {
		\fpSet \lTmpbFp {\fpEval{\cellGetText{##1}{#1}}}  
		\fpSet \lTmpcFp {\fpMathDiv {\lTmpbFp} {\lTmpaFp}}
		\tlPutRight \gTikzPercentagePathsTl {
			\evalWhole{
				(##1-#1.south~west) rectangle
				($(##1-#1.north~west)!{\fpUse \lTmpcFp}!(##1-#1.north~east)$)
			}
		}
	}
}
\IgnoreSpacesOff
